% Part: lambda-calculus
% Chapter: lambda-definability
% Section: introduction

\documentclass[../../../include/open-logic-section]{subfiles}

\begin{document}

\olfileid{lam}{ldf}{int}
\olsection{مقدمة}

من نظر إلى حساب لامبدا أول الأمر رآه حسابًا مجردًا للعبارات التي تمثل
الدوال وتطبيق بعضها في بعض. وليس في تركيبه ما يخص هذه الدوال بنوع
معين من الأشياء؛ بل لا يذكر هذا التركيب الأعداد الطبيعية أصلًا.
فالتطبيق وتجريد لامبدا، وهما عمليتاه الأساسيتان، يجريان على الدوال
كيفما كانت، ولا يختصان بالدوال على الأعداد الطبيعية.

ومع هذا العموم نستطيع، بحسن التصرف في الترميز، أن نعرّف فيه الدوال
الحسابية، أي الدوال على الأعداد الطبيعية. وسبيل ذلك أن نخص كل عدد طبيعي
$n \in \Nat$ بحد~$\lambd$ نسميه~$\num n$، وهو \emph{عدد تشيرش}
الذي يمثل~$n$. وهذه الأعداد منسوبة إلى ألونزو تشيرش.

\begin{defn}
  لكل~$n \in \Nat$ \emph{عدد تشيرش}~$\num{n}$ الذي يمثل~$n$،
  وتعريفه
  \[
    \num{n} \ident \lambd[fx][f^n(x)]
  \]
  والمراد بـ$f^n(x)$ حاصل تكرار تطبيق~$f$ في~$x$ بمقدار~$n$ من المرات.
  فـ$\num{0}$، مثلًا، هو~$\lambd[fx][x]$، و$\num{3}$
  هو~$\lambd[fx][f(f(f\,x))]$.
\end{defn}
  
فترميز عدد تشيرش~$\num n$ حد لامبدا يمثل دالة ذات وسيطين~$f$ و$x$،
تردّ القيمة~$f^n(x)$. وهذه الأعداد، كما يظهر من تعريفها، سوية.

وليس المقصود مجرد تمثيل الأعداد الطبيعية، بل التمكن من الحساب بها.
فإذا حسبنا بأعداد تشيرش في حساب لامبدا، طبقنا حدًا من حدود~$\lambd$،
وليكن~$F$، في عدد تشيرش، ثم اختزلنا الحد المركب~$F\, \num n$ إلى صورة سوية.
فإن انتهى هذا الاختزال دائمًا إلى صورة سوية، وكانت دائمًا عدد تشيرش~$\num m$،
عددنا نتيجة الحساب هي العدد~$m$.
وبهذا يكون~$F$ تعريفًا لدالة~$f\colon \Nat \to \Nat$،
تتحقق فيها~$f(n) = m$ إذا وفقط إذا~$F\, \num n \red \num m$.
ووحدانية الصورة السوية، إن وجدت، مضمونة بخاصية تشيرش--روسر؛
فإذا تحقق~$F\, \num n \red \num m$، امتنع أن يكون ثمة حد سوي آخر،
وبالأخص عدد تشيرش آخر، يختزل إليه~$F \, \num n$.

ولنعكس النظر: إذا أعطينا دالة~$f\colon \Nat \to \Nat$،
فهل نجد حدًا~$F$ يعرّف~$f$ على هذا الوجه؟
إن وجد قلنا إن~$F$ \emph{!!{lambda define}s}~$f$،
وإن~$f$ !!{lambda definable}. ثم نمد هذا الاصطلاح إلى الدوال
المتعددة المواضع وإلى الدوال الجزئية.

\begin{defn}
لتكن~$f\colon \Nat^k \pto \Nat$ دالة جزئية ذات~$k$ مواضع.
شرط أن يكون حد لامبدا~$F$ \emph{!!{lambda define}s}~$f$ هو أنه،
لكل~$n_0$، \dots،~$n_{k-1}$،
\[
F \, \num{n_0} \, \num{n_1} \dots \num{n_{k-1}} \red \num{f(n_0, n_1, \dots,
  n_{k-1})}
\]
إذا كانت~$f(n_0, \dots, n_{k-1})$ معرّفة؛ وأما إذا لم تكن معرّفة،
فلا صورة سوية للحد~$F \, \num{n_0} \, \num{n_1} \dots \num{n_{k-1}}$.
\end{defn}

وأيسر أمثلة ذلك الدوال الثابتة: فالحد~$C_k \ident
\lambd[x][\num{k}]$ !!{lambda define}s الدالة~$c_k\colon \Nat \to
\Nat$ التي تعطي~$c_k(n) = k$، لأن
$C_k \, \num n \ident (\lambd[x][\num{k}])\num n \redone \num{k}$
لكل~$n$. وأما الدالة المحايدة فهي !!{lambda defined} بالحد~$\lambd[x][x]$.
وكلما زاد تركيب الدالة عسر تعريفها، وربما احتاج إلى تصرف دقيق.
فلعل مما يدعو إلى العجب أن كل دالة قابلة للحساب هي
!!{lambda definable}. ويصح عكس ذلك أيضًا:
فكل دالة !!{lambda definable} قابلة للحساب.

\end{document}
